\documentclass[a4paper]{book}
\usepackage{amsmath}
\usepackage{ctex}
\usepackage{fancyhdr}
\usepackage{titlesec}
\usepackage{amssymb}
\usepackage[margin=1in]{geometry}
\usepackage{graphicx}   % 插入图片支持
\title{\fangsong\zihao{1}逆光而行，与子同途}
\author{$\hbar $i1ber$\dag$}
\begin{document}
\maketitle
    
\newpage

\setlength{\parindent}{2em}
\setlength{\parskip}{0.5em}
\pagestyle{fancy}
\fancyhf{}
\fancyhead[LE]{\footnotesize \thepage} % 偶数页(LE)左侧：当前章节名
\fancyhead[RO]{\footnotesize \thepage} % 偶数页(LE)左侧：当前章节名
\fancyhead[CO]{\footnotesize 逆光而行，与子同途} % 奇数页左侧(LO)和偶数页右侧(RE)：书名
\fancyhead[CE]{\footnotesize \leftmark} % 奇数页左侧(LO)和偶数页右侧(RE)：作者名
\setcounter{chapter}{-1}
\begin{center}
    \songti\zihao{-2}
    这本书由 \heiti\zihao{-2} {攻打赫尔利！} 出版。\\
    \end{center}
    \heiti\zihao{3}{故事简介}\\
    \begin{figure}[htbp]
    \centering
    \includegraphics[width=0.5\textwidth]{111.png} % 替换为实际图片文件名
\end{figure}
    

    \songti\zihao{4}
在崇尚绩点与排名的K大，789和278是典型的“差生”——一个国护队标兵GPA2.9，一个后勤骨干GPA2.5。他们是彼此最硬的铠甲，也是最深的软肋，在训练场的汗水中淬炼出“死生契阔”的默契，却也在学业的自卑与醋意中暗生裂痕。

当789因学业焦虑而失控，当278在生死关头被789舍身相护，一句“我答应过要亲手给你戴上学士帽”的誓言，成为了照亮绝境的曙光。然而，现实的重量如期而至：企业的橄榄枝与家庭的召唤，仿佛要将他们推向不同的人生轨道。

这是一段偏离主流赛道的青春。

在这里，荣耀不只在成绩单上，更在每一次并肩冲锋的背影里；

在这里，爱情没有甜言蜜语，只有在绝境中脱口而出的承诺，和深夜阳台上一句“你去哪里，我的后勤就跟到哪里”的坚定选择。

且看两个光着膀子在宿舍楼晃荡的少年，如何逆着世俗的偏见，踏过荆棘与迷茫，用他们独有的方式，定义属于自己的巅峰，走向那盏名为“彼此”的灯火归处。

这是一个关于勇气、选择和爱的故事。它告诉我们，真正的光明坦途，是与那个对的人，携手逆光而行，走出属于自己的风景。

\tableofcontents

\large
\chapter{料峭春风中的烂漫}
\songti\zihao{4}
\setcounter{page}{1}

清明前的风，还带着料峭的寒意，却已吹醒了K大校园里那几株百年海棠。粉白的花瓣如云如霞，簌簌地落在一行身着礼宾服、步伐铿锵的年轻人肩头。

K大国护队承担了此次清明祭扫英烈活动的仪仗任务。队伍最前方，是掌旗手兼刀手——789。他戴着白手套，紧握着冰冷的仪仗指挥刀，每一步都踏得沉稳如山，刀锋在微弱的晨光中划出冷冽的弧线。他神情肃穆，目光平视远方，仿佛周遭的海棠花雨、窃窃私语的人群，都已与他无关。他是整个仪仗的灵魂，承载着无比的庄重与荣耀。

活动吸引了很多人观摩，除了师生，还有附近小学组织来的孩子。孩子们天性活泼，在肃穆的仪式间隙，目光很快被国护队中另一个身影吸引——负责护旗的278。

278的身高和挺拔的身姿在队列中本就出众，加之他长相帅气，即使穿着统一的礼服，也难掩那份阳光俊朗。仪式刚结束，还没等完全解散，几个胆子大的小朋友就呼啦啦围了上去，仰着红扑扑的小脸，叽叽喳喳地要求合影。

“哥哥你好帅！”

“哥哥，我能摸摸你的绶带吗？”

面对突然涌来的“崇拜者”，278有点措手不及，平时的散漫不羁收了起来，脸上露出一丝难得的、混合着不好意思和欢喜的腼腆。他蹲下身，配合着孩子们，嘴角不自觉地扬起，笑容干净又明亮。

这时，校报记者也趁机上前采访：“同学，请问作为国护队成员，参加这次清明祭扫活动，你有什么感想？”

278抱着一个扑到他怀里的小男孩，闻言愣了一下，似乎没准备措辞。他张了张嘴，话在喉咙里打了个转，却只冒出几个零碎的词：“啊？感想……就是，嗯……挺光荣的，对……光荣……” 语言贫乏得有些“个楞”，与他刚才和孩子们互动时的自然判若两人。

就在278卡壳，脸上开始泛起尴尬的红晕时，一道平静无波的目光落在他身上。是刚刚完成刀手任务、正肃立在一旁的789。他没有说话，只是极轻微地、几不可察地朝278点了点头，眼神里没有责备，只有一种沉静的力量，仿佛在说：“慢慢想，说心里话。”

接触到789的目光，278像是突然找到了主心骨，慌乱的心瞬间安定下来。他深吸一口气，重新看向记者，语气依旧不算流畅，却多了几分真诚：“就是……觉得不能忘记吧。我们现在的安稳日子，是前辈用命换来的。穿着这身衣服站在这里，……觉得责任挺重的。”

他的回答朴实无华，却比任何华丽的辞藻都更能打动人心。记者满意地记录着。789的目光依旧停留在278侧脸，那眼神深处，掠过一丝极淡的、连他自己都未曾察觉的柔和。

不远处，一树繁茂的海棠下，A静静地站着。她也来观摩活动，但更像一个冷静的旁观者。她没有看被孩子和记者围住的278，也没有看肃立如松的789，她的目光越过喧闹的人群，落在远处纪念碑上苍劲的文字。

当哀乐响起，全场默哀时，她微微垂着眼睑，嘴唇无声地翕动，背诵的是屈原的《九歌·国殇》。当念到最后一句时，她的声音几不可闻，却带着一种刻意强调的坚定：

“身既死兮神以灵，魂魄毅兮为鬼雄。”

她希望此刻充盈自己心间的，是这磅礴悲壮的家国情怀，是对于牺牲与奉献的追思。她用力地想着，试图用这宏大的意念，去覆盖、去驱散心底那抹两年前因某个“并非第一梯队”的天才少年幻灭而留下的、对于“不够完美”的失望与阴影。

海棠花瓣依旧无声飘落，落在少年们挺直的肩背，落在少女微颤的睫毛上。寒食时节的空气里，弥漫着花香、哀思、荣光，以及悄然滋长、尚不自知的情愫与坚持。

这是一个关于铭记、关于成长、关于在绚烂与肃穆之间，寻找自身坐标的故事的开端。一切，都始于这个海棠烂漫的时节。
\newpage
\large
\chapter{深埋的默契}
\normalsize
\section{两种“差生”}
\songti\zihao{4}
初夏的午后，阳光透过《电路电磁学》教室高大的窗户，在弥漫着粉笔灰和沉闷空气的房间里，切出明晃晃的光带。讲台上，老师正用毫无起伏的语调念着期中考试的分数和学号。

“789，65分。”

坐在后排的789脊背微微一僵。他伸手拿回卷子，目光扫过那些刺眼的红叉，指关节下意识地收紧，在卷面上留下几道浅痕。65分。这个分数像一根针，扎在他因为前一晚加练而依旧酸胀的肌肉上。他不自觉地挺直了腰板，那是长期站军姿形成的本能，仿佛这样就能撑起那点微不足道的尊严。对他而言，这个分数不是“及格”，而是“本该更好”的失败。国护队训练占用的每一分钟，都在这个数字面前变成了沉甸甸的代价。

“278，40分。”

不远处，他的哥们儿278闻言，几乎是蹦跳着上去拿了卷子，脸上看不出半点羞愧，反而带着一种“终于落地”的坦然。他回到座位，把卷子往桌上一拍，发出“啪”的一声轻响，引来旁边几个同学侧目。他浑不在意，甚至用笔戳了戳前排一个男生的后背，小声调侃道：“嘿，比你高两分，今晚夜宵你请。”

“A，95分。”

一个坐在前排靠窗位置的女生应声站起来。她身形高挑，穿着简单的T恤牛仔裤，短发清爽，侧脸线条在光线下显得有些模糊不清。她平静地取回卷子，脸上没有任何喜悦的表情，仿佛这只是完成了一个微不足道的流程。她是物理系的A，这门跨专业选修课对她来说似乎游刃有余。

下课铃响，教室里瞬间嘈杂起来。278抓了抓头发，目光在教室里逡巡，最后落在了正在收拾书包的A身上。他像是下了很大决心，抓起那张40分的卷子，挤过人群，拦在了对“死刷题的木头脑袋”深恶痛疾的A面前。

“那个……A同学，”278努力让自己的笑容显得真诚无害，“这道大题，你能给我讲讲不？答案我看不懂。”

A抬起头，目光平静地扫过278，又落在他手指着的那道题上。她没有接卷子，只是淡淡地开口，声音清晰却带着一种疏离的冷感：“这道题考察的是高斯定理在非对称场中的应用。但你连最基础的静电场环路定理前提都没搞清，第二小问的公式用得完全是错的。”她抬起眼，直视着278有些窘迫的脸，“这几个公式是前置基础，你没弄懂之前问这个没意义。”

她的评价像手术刀一样精准，瞬间剖开了278试图蒙混过关的侥幸心理。278脸上的笑容僵住了，他讪讪地收回卷子，嘟囔了一句“哦，谢谢啊”，便转身败下阵来。A的话不多，却像一堵无形的墙，让他彻底放弃了在学业上向她求助的念头。

挫败感只持续了几秒钟。278一扭头，正好看见同班的C也从讲台上拿回了卷子——75分。C是班上的活跃分子，成绩中上，为人热心。278立刻像找到了救星，脸上瞬间阴转晴，几步凑了过去。

“C哥！牛逼啊！75分！”他亲热地揽住C的肩膀，“最后那道题怎么回事？快给我讲讲，哥们儿今晚能不能睡个好觉就靠你了！”

C笑着推了他一下，显然很吃这一套。“你小子，又没复习吧？走，去自习室，正好我也要整理一下笔记。”

两人勾肩搭背，有说有笑地朝教室外走去。C拿出自己条理清晰的笔记本，278凑在旁边，一边听一边夸张地“哦~”、“原来如此！”地附和着，气氛融洽得刺眼。

后排，789将这一切尽收眼底。

他看着278在A那里碰了一鼻子灰，又看着278转而和C相谈甚欢。他的目光落在278自然地搭在C肩上的手臂，落在C那本干净整齐的笔记上，最后落回自己手里这张皱巴巴的、只有65分的卷子。

一股莫名的烦闷，像窗外的热浪一样，毫无征兆地涌上心头。比刚才拿到分数时更甚。

他独自坐在渐渐空荡的教室里，阳光把他的影子拉得很长。训练场上的汗水和荣耀，在此刻被这间闷热的教室和几个冰冷的数字衬得有些遥远。他深吸一口气，将卷子胡乱塞进背包，起身离开。坚硬的军靴底敲击着地面，发出沉闷的回响，与窗外278和C渐行渐远的笑闹声，格格不入。
\normalsize
\section{战场与软肋}
\songti\zihao{4}
夕阳将天边烧成一片绚烂的橘红，但K大东操场的塑胶跑道上，热度并未完全散去，混合着青草和尘土的气息，蒸腾起一股独属于训练场的、充满力量感的热浪。

这里，是另一个世界。是789和278的战场。

“正步——走！”

口令声如同出鞘的利剑，划破空气。队列最前方，789的身影如同标枪般挺拔。他的每一个动作都像是用刻度尺量过，手臂摆动的角度，腿踢出的高度，甚至脚面下压的瞬间，都带着一种刀劈斧凿般的力量感和精确性。汗水从他剃得极短的发茬间渗出，沿着晒成小麦色的脖颈滑落，洇湿了草绿色的作训服后背，但他眼神锐利，神情专注，仿佛周遭的一切都已消失，只剩下节奏、力量和不容置疑的标准。

他是国护队的标兵，是队伍的门面，是训练场上当之无愧的焦点。

而在场地边缘，树荫笼罩下的器材区和休息区，则是278的领地。他没那么挺拔如松，反而微微弓着腰，面前摊开一个巨大的军用迷彩背包，里面分门别类地装着水壶、药品、备用零件、擦枪工具。他手里正飞快地检查着一副护膝的搭扣，眼神专注，动作麻利，带着一种不同于课堂上的认真。

当队伍在口令下暂时休息，队员们或活动手脚，或冲向场边喝水时，278的工作才真正开始。他像一只敏捷的土拨鼠，在人群中穿梭，递水，检查装备，处理一些突发的小状况。

789没有像其他人一样立刻放松。他走到场边，习惯性地在一个马扎上坐下，开始解自己膝盖上那副已经有些磨损的护膝。长时间的踢腿练习让他的膝盖旧伤有些隐隐作痛，这是他作为标兵必须付出的代价。

278晃悠了过来，看似随意地拿起789刚解下的护膝，掂量了一下，眉头微皱：“这副快不行了，卡扣有点松，里面的海绵也压薄了。” 他一边说着，一边从自己的百宝袋里熟练地翻出一副看似一模一样、实则崭新的护膝，递了过去。“先用这副，我晚上把旧的修修。”

789“嗯”了一声，接过来，没有任何多余的言语，直接开始佩戴。然而，当他的手指探入护膝内侧时，动作微微一顿。他敏锐地感觉到，在膝盖正前方的位置，里面似乎多了一层异常柔软且有弹性的东西。他抬头看了278一眼。

278正拧开自己的水壶猛灌了几口，感受到789的目光，他抹了把嘴，咧嘴一笑，露出一口白牙，压低声音：“加了点儿好东西，高分子缓冲凝胶，搞竞赛那帮人用的，舒服吧？”

那是278不知道从什么渠道弄来，又偷偷缝进去的软垫。这小子总有这些稀奇古怪的门路，而且总能用在他认为值得的地方。

789没说话，只是低下头，继续系紧卡扣。那多出来的一层柔软，恰到好处地包裹住他酸胀的膝关节，带来一种细微却真实的慰藉。一种无需言谢、早已融入骨血的习惯与默契。

训练继续，强度升级。当最后一组冲刺跑结束，夜幕已悄然降临。789几乎是拖着步子走到场边，汗水像雨一样往下淌，胸膛剧烈起伏，整个人像是从水里捞出来。他连说话的力气都没有，直接瘫坐在了地上。

278的情况好不少，他主要负责后勤保障，体力消耗没那么大。他拿着两瓶功能饮料走过来，先拧开一瓶，自己咕咚咕咚喝掉大半，然后很自然地把剩下的小半瓶递到789面前。

“喏，给你留了口。”

789抬眼，汗水模糊了他的视线，他只看到278递过来的那个瓶子，和瓶口那点不多的水。他没有丝毫犹豫，甚至没抬手去接，只是就着278的手，微微仰头，几口就将那点水喝得一滴不剩。喉咙里火烧火燎的感觉稍稍缓解。

这个动作做得如此理所当然，仿佛已经重复过千百遍。分享同一瓶水，在他们是比课堂上传纸条更寻常的事。

278收回空瓶子，看着789疲惫却放松的侧脸，用肩膀轻轻撞了他一下，语气带着点戏谑，却又藏着不易察觉的关切：“还行不行啊，标兵？别明天趴窝了，我可背不动你。”

789终于缓过一口气，瞥了他一眼，声音因脱力而沙哑，却带着训练场上独有的硬气：“趴不了……还得给你挣夜宵呢。”

夜色中，操场上的灯光亮起，将两人的影子拉长，交融在一起。课堂上的分数和挫败，仿佛已是上个世纪的事。在这里，他们是彼此最坚硬的盔甲，也是最清楚的软肋。这种在汗水中淬炼出的羁绊，远比一张满分的试卷，更来得牢固和真实。
\normalsize
\section{阳台上的星与尘}
\songti\zihao{4}
夜很深了，白天的燥热褪去，只剩下清凉的晚风，裹挟着夏夜草木的湿气，轻轻拂过皮肤。宿舍楼大多已经熄灯，只有零星几个窗口还亮着，像是沉睡巨兽未合的眼。

789和278并排靠在三楼阳台冰凉的铁栏杆上，两人都只穿着宽松的篮球裤衩，光着上身。白天的汗水和尘土已经被冷水冲刷干净，皮肤在朦胧的夜色里泛着微光。789的身材确实练得极好，胸肌腹肌块垒分明，在放松状态下也勾勒出清晰的线条，那是长期高强度训练刻下的勋章。278则显得更“家常”些，高大挺拔，骨架匀称，身上覆盖着一层结实但不过分贲张的肌肉，是那种长期活动而非刻意雕琢出的健壮。

空气里飘着淡淡的沐浴露香和一丝若有若无的膏药味——789的膝盖旧伤在今晚的训练后又有些发作，278刚帮他重新贴了活血止痛的膏药。

周围很静，只有不知名的虫鸣和远处马路隐约传来的车流声，像是这个世界沉睡时的呼吸。

789仰着头，望着被城市光污染映得有些发灰、却依旧有零星几点亮光顽强闪烁的夜空。他很久没说话，只是静静地望着，下颌线绷得有些紧。白天训练场上的锐利和坚定，在此刻仿佛被夜色稀释，透出一种难得的、近乎迷茫的疲惫。

“有时候觉得，”他忽然开口，声音比平时低沉沙哑不少，融在夜风里，几乎听不真切，“离开操场，离开国护队，我……好像什么都不是。”

这句话很轻，却像颗小石子，投入了安静的夜色里。

278正叼着根没点燃的烟过干瘾，闻言，侧过头看他。月光勾勒出789硬朗的侧脸，但那眼神里一闪而过的空虚，没能逃过278的眼睛。他太熟悉789了，熟悉他每一个表情背后的情绪。

278没立刻接话，只是用肩膀不轻不重地撞了一下789的肩膀，力道刚好让他从那种自我放逐的情绪里晃神过来。然后，他咧开嘴，露出一个在夜色里依旧晃眼的白牙笑容，语气带着他特有的、混不吝的调侃：

“胡说八道什么呢你。”

他把烟从嘴上拿下来，夹在手指间虚点着789：“离开那里，你789还是我278的爹！亲爹！期末考还得指望你划重点、传答案、带我悬梁刺股闯过及格线呢，懂不懂？”

这话糙理不糙，甚至有点混账，但奇异地驱散了那点萦绕在789周身的低气压。他把789那点关于存在价值的哲学思辨，一下子拉低到了“哥们儿需要你救命”的世俗层面，反而显得无比真实和接地气。

789被他这话噎了一下，下意识扭头瞪他，却对上278笑得没心没肺却异常温暖的眼神。那眼神里明明白白写着：别想那些没用的，咱俩谁跟谁，你牛逼你的，我抱我的大腿，这不天经地义？

看着278那副理所当然、“离了你我活不了”的赖皮样子，789心里那点突如其来的脆弱和虚无，忽然就散了。他嗤笑一声，带着点无奈，也带着点释然，抬手给了278结实的胳膊一拳：“滚蛋，谁是你爹。”

“你啊！”278揉着胳膊，笑得更大声了，引得楼下似乎传来一声不满的咳嗽。他赶紧压低声音，凑近些，用气声道：“说真的，别瞎想。你就是你，甭管在哪儿。在我这儿，你永远是扛把子那个。”

夜色掩去了789脸上可能泛起的细微波动。他没再说话，只是重新靠回栏杆上，再次抬起头看向夜空。但这一次，他的肩膀不再紧绷，呼吸也变得悠长而平稳。

星空依旧稀疏，尘世依旧喧嚣。但在这个小小的阳台上，两个只穿着裤衩的年轻人，靠着冰冷的栏杆，用他们之间独有的、粗糙又直接的方式，完成了一次无声的充电和更深的绑定。他们是彼此在迷茫夜色里，最能触摸得到的真实和依靠。
\normalsize
\section{醋意的裂痕}
\songti\zihao{4}
几天后的晚上，789想去图书馆找点资料。穿过灯火通明的自习区时，他的脚步不自觉地慢了下来。在密集的书架和埋头苦读的人群中，他几乎一眼就看到了那个熟悉的身影——278。

278又和C坐在一起。这次不是在讨论题目，C正眉飞色舞地说着什么，278则侧着头，听得津津有味，脸上挂着789很久没见到过的、带着点好奇和八卦的轻松笑容。

789本想径直走开，但C提高了些许的音调，还是断断续续地飘进了他的耳朵：

“……所以说，别看A现在一副‘凡人勿近’的样子，高中时候也挺……啧，你猜怎么着？”C压低了声音，带着分享秘密的兴奋，“她那时候好像挺崇拜她们学校一个搞物竞的学长，那学长计算机、数学物理都挺牛，在他们那小地方算是顶尖了，不过放到省里嘛……嗯，不算最拔尖那拨。”

278听得入神，催促道：“然后呢？这跟电路有什么关系？”

“哎，听我说嘛！”C继续道，“就那次电路竞赛培训前，不是快到六一了嘛？不知谁起哄，说给学长送个儿童节礼物祝他‘永葆童心’、竞赛顺利。A她……她居然真的准备了，好像是自己焊的一个什么小电路板，据说还挺精巧。结果送去的时候被人撞见，好一阵起哄，搞得她后来看见那学长都绕道走……哈哈，没想到吧？”

278也跟着笑了起来，摇着头感叹：“真没看出来……不过她那种靠灵性学习的，会喜欢上同样有天赋但又不是‘刷题机器’的学长，好像也挺合理？”

“可不是嘛！”C附和道，“所以啊，这种人……”

后面的话，789没再听清。他只觉得一股无名火“噌”地窜了上来，烧得他胸口发闷。A的往事，278听得那么投入，还和C一起讨论、一起笑？那种融洽的氛围，那种分享秘密的亲近感，像一根细小的刺，扎得他极不舒服。他猛地转身，离开了自习区，连原本要找什么资料都忘了。

晚上九点，本该是休息的时间，东操场却再次响起了沉重而单调的脚步声。

只有789一个人。他没有开灯，借着远处路灯和月光，在跑道上进行着障碍穿越加练。他的动作比白天更猛、更急，每一次跳跃、翻滚、匍匐，都带着一股要把自己拆解开的狠劲。汗水早已浸透作训服，沉重的喘息声在寂静的夜里格外清晰。这已经不是训练，更像是一种惩罚，一种发泄。

就在这时，278的身影出现在跑道边缘。他大概是听说了789又来自虐加练，不放心地找了过来。看着789那近乎疯狂的样子，278皱紧了眉，快步走了过去。

“789！你他妈疯了吗？快停下！”278试图拦住他。

789猛地停下动作，撑着膝盖，大口喘着气，抬起头。汗水顺着他的下颌线滴落，眼神在夜色里烧得通红。他盯着278，胸口剧烈起伏，那句憋了一晚上的话，终于失控地冲口而出：

“不去找你的C学长了？不去听他讲A的八卦了？”

话音落下，空气瞬间凝固。278愣住了，脸上的担忧和急切僵在那里，转而变成了一种难以置信和受伤的表情。他张了张嘴，似乎想解释什么，但看着789那副浑身是刺、拒人千里的样子，最终什么也没说，只是眼神一点点暗了下去。

也就在这一刻，操场入口处，一个高挑的身影短暂停留。是A。她似乎刚从实验室回来，手里还拿着几本书。她看到了跑道中央对峙的两人，看到了789那副自虐般的状态和两人之间紧绷的气氛。她的目光在789身上停留了两秒，眉头微不可察地蹙了一下，那双习惯于分析公式和定理的眼睛里，清晰地闪过一丝不解。

在她看来，这种毫无效率、纯粹消耗体能和健康的情绪化发泄，是极其不理智且毫无意义的。她无法理解这种被原始情绪驱动、近乎野蛮的行为逻辑。她没有停留，甚至没有改变步速，就像路过一个无法理解的噪音源，很快便消失在操场的另一个出口，仿佛从未出现过。

夜色重新笼罩下来，比刚才更沉、更冷。789和278站在原地，中间隔着的几步距离，却仿佛突然变成了一道裂开的深渊。第一次，充满默契和汗水的训练场，只剩下无声的僵持和冰冷的隔阂。

\newpage
\large
\chapter{惊雷与誓言}
\normalsize
\section{联合集训}
\songti\zihao{4}
省高校军事技能联合集训基地，坐落在远离市区的山脚下。这里的气氛和K大校园截然不同，更加肃杀，也更加直接。来自全省十余所高校的国护队、预备役标兵们齐聚于此，空气中弥漫着一种无声的竞争意味。

789和278穿着统一的K大作战训服，站在本校的队伍里。789身姿如松，眼神锐利地扫视着其他学校的队伍，像一头进入新领地的豹子，充满了警惕和征服欲。278则显得放松些，但那双平时总带着笑意的眼睛此刻也认真了不少，同样在打量着潜在的对手，尤其是他们的装备。

开营仪式后的第一次混合编队体能拉练，就成了没有硝烟的战场。各校队员混编成几个临时小队，进行负重越野。789所在的小队里，有几个来自以学风严谨、学生成绩优异著称的H大的队员。

途中休息时，大家卸下背囊，喘着气喝水。一个H大的队员，看着臂章上是“机电工程”专业，扶了扶眼镜，语气带着几分优越感，看似随意地提起：“K大的兄弟，你们学校环境是真好，听说社团活动特别丰富？就是学业压力是不是相对……轻松点？”

这话听起来是闲聊，但“轻松点”三个字，配上那意有所指的眼神，分明是在影射K大学生普遍“爱玩不学习”的刻板印象。另一个H大的也笑着接话：“是啊，看你们国护队搞这么大阵仗，平时训练肯定占不少时间吧？GPA还能顾得上吗？佩服佩服。”

这种看似客气实则带刺的话，让K大另外几个队员脸色都有些难看。278正拧开水壶，闻言，脸上的懒散瞬间收了起来，他嗤笑一声，声音不大但足够清晰：“哟，H大的兄弟就是不一样，跑个步都惦记着GPA。怎么，是怕这次集训评分不够，影响你们保研加分？”

他这话直接又刁钻，一下子把对方那点弯弯绕捅破了。那几个H大队员脸色顿时有些挂不住。

789一直没说话，但下颌线绷得死紧。H大队员的话精准地戳到了他内心最敏感的痛处——学业。他因为国护队训练付出的时间、精力，以及那始终不上不下的GPA，在此刻成了别人轻蔑的谈资。一股混合着屈辱和愤怒的火直冲头顶。

他猛地站起身，背囊被他拎起时发出沉闷的响声。他盯着刚才说话的那两个H大队员，眼神冷得像冰。

“废话那么多，是靠嘴皮子拉练吗？”789的声音不高，却带着一股压不住的狠劲，“是骡子是马，拉出来遛遛。后面还有障碍赛和战术协同，到时候再看，谁的训练是花架子，谁的成绩是纸上谈兵！”

他的反应极其强硬，甚至有些过激，一下子把矛盾摆上了台面。周围瞬间安静下来，其他学校的队员都看了过来，气氛变得剑拔弩张。

278见状，也立刻站了起来，不动声色地站到789侧前方半步，看似随意，却是一个隐隐维护的姿态。他脸上又挂起了那种混不吝的笑，打圆场道：“行了行了，都一个战壕的兄弟，较这劲干嘛。我们K大别的不行，就是体力好，不服气一会儿赛场上见真章呗？”

他这话既接了789的挑衅，又把焦点引回了训练本身，给了对方一个台阶，也避免了冲突进一步升级。

H大的队员哼了一声，没再说话，但眼神里的轻视并未减少。休息结束的哨声响起，队伍重新出发。

接下来的路程，789像是跟谁赌气一样，冲得特别猛，完全超出了常规的节奏。他的每一个动作都充满了爆发力，仿佛要把所有被轻视的怒火都发泄在跑道上。278跟在他身后，眉头微蹙，眼神里透出担忧。他知道789的老毛病又犯了——一遇到和学业、和外界评价相关的刺激，就容易变得激进，用近乎自虐的训练来证明自己。

这种状态，在竞争激烈的集训中，就像一把双刃剑。

278加快脚步，凑到789身边，压低声音：“慢点，蓄点力！后面项目还长着呢！”

789却像没听见，反而又加快了几分速度，只留给278一个紧绷而倔强的背影。

山风吹过林间，带来一丝凉意，却吹不散789心头那团灼烧的火焰，也吹不散278眼中越来越浓的忧虑。裂痕尚未修补，外部的压力已然袭来。这场集训，注定不会平静。
\normalsize
\section{意外降临}  
\songti\zihao{4}
联合集训的第四天，迎来了最严酷的挑战：夜间野外负重拉练与障碍穿越。路线设置在基地后方未经深度开发的山区，地形复杂，充满了未知。

经过几天高强度的训练和内心的煎熬，789的体力消耗已接近临界点，但他的精神却因一股不服输的狠劲而异常亢奋。在H大队员若有若无的注视下，他几乎是下意识地冲在K大队伍的最前面，仿佛要用这种方式碾碎所有质疑。

“789，慢点！这段路况不好！” 278在他身后不远处喊道，声音带着明显的焦急。夜色浓重，仅靠头灯的光束难以完全看清脚下的每一处细节。

789像是没听见，或者说，他听见了，但体内那股急于证明什么的火燃烧得太旺，压过了理智。他加快脚步，试图率先通过一段需要借助固定绳索攀下的陡坡。

就在这时，意外发生了。

就在789下行、278紧随其后，双手刚完全承重于那根主绳索时，上方一个原本应该牢固嵌入岩体的保护点，可能因为长期风化承重，也可能是之前队伍使用过度，竟突然发出了令人牙酸的“嘎吱”声，随即猛地松脱！

“小心！” 789在下方看得真切，瞳孔骤缩，嘶声大吼。

但警告来得太迟。278只觉得手上一轻，整个人瞬间失去平衡，伴随着碎石滚落的声音，猛地向下坠去！

那一刻，时间仿佛被拉长。278的头灯光线在黑暗中疯狂乱晃，失重的恐惧瞬间攫住了他。

“278！”

789的反应快得超出了思考。几乎是在保护点松脱的同一瞬间，他根本没有权衡利弊，身体已经先于大脑做出了行动。他非但没有躲闪，反而逆着278坠落的方向，用尽全力猛地向上扑了过去，不是去抓那根已然失效的绳子，而是用自己的身体，悍然撞向失控下坠的278！

他试图将278撞向内侧的岩壁，争取缓冲和抓住其他固定物的机会。

“砰！”

沉重的撞击声闷闷地响起。789成功地改变了278的下坠轨迹，两人的身体狠狠撞在陡坡上，又一起向下滑落了七八米，才被一堆茂密的灌木和一块凸起的岩石险险卡住。

短暂的死寂后，是278带着哭腔的、颤抖的呼喊：“789？！789！你怎么样？！”

278发现自己除了多处擦伤和惊吓，并无大碍。但他身下的789，却在他动弹的瞬间，发出了一声压抑不住的、极其痛苦的闷哼。

278慌忙翻身，借助滚落到一旁的头灯光线，看清了眼前的景象，血液几乎瞬间冻结——

789的脸色惨白如纸，额头上全是冷汗，他的左臂以一个不自然的角度弯曲着，显然是脱臼或骨折了。而最触目惊心的是，他的右边小腿被一段在撞击中断裂的、尖锐的岩石棱角划开了一道深可见骨的口子，鲜血正汩汩地涌出，迅速染红了作训服和身下的泥土。

“789！你……” 278的声音彻底变了调，手忙脚乱地想要按住伤口，却不知从何下手，眼泪不受控制地涌了出来。是因为疼痛，更是因为后怕和滔天的悔恨。如果不是为了救他，789根本不会……

“别……哭……” 789的声音极其虚弱，却带着一种奇异的镇定。他抬起没受伤的右手，艰难地碰了碰278满是泪水和污泥的脸，“死……死不了。”

寒冷和失血让他开始控制不住地发抖。278脱下单薄的作训服外套，徒劳地想给他止血、保暖，自己则只剩下了一件湿透的短袖，在夜风里冷得牙齿打颤。他一边用尽全力按压着789腿上的伤口，一边朝着漆黑的夜空，用尽全身力气嘶吼：“有人吗？！救命——！这里有人受伤了——！”

回应他的，只有山风的呼啸和远处隐约的、其他小队行进的声音。绝望，像冰冷的潮水，一点点淹没上来。

278紧紧抱着怀里不断失温、发抖的789，巨大的恐惧和无力感几乎要将他击垮。他语无伦次地哽咽着：“对不起……对不起……都是我不好……你坚持住……一定会有人来的……”

在逐渐模糊的意识中，789能感觉到278滚烫的眼泪滴落在自己脸上，能听到他心脏因为恐惧而疯狂的跳动。一股从未有过的、比身体上的疼痛更尖锐的情绪攫住了他。他用尽最后一丝力气，抬起右手，死死攥住了278按压在他伤口上的那只手。

手指冰冷，却带着不容置疑的力量。

“听着……” 789的声音气若游丝，却异常清晰，每一个字都像是从胸腔里挤出来的承诺，“别怕……我们……都要活下去……”

他顿了顿，涣散的目光努力聚焦在278泪眼模糊的脸上，仿佛要穿透这绝望的黑暗，看到某个确定的未来。

“我答应过……”他几乎是咬着牙，许下了这个在绝境中如同灯塔般的誓言，“要……亲手……给你戴上……学士帽……”

这句话说完，他最后一点力气也耗尽了，脑袋一歪，彻底昏厥在278的怀里。

“789——！” 278的哭喊声撕裂了寂静的山谷。他死死抱着怀里失去意识的人，像一头被困的幼兽，在无边无际的黑暗和寒冷中，唯一能感受到的温暖和支撑，只剩下彼此紧握的手，和那个在生死边缘，用尽全力许下的、关于未来的约定。
\normalsize
\section{无声的确认}
\songti\zihao{4}
意识先于视觉回归。首先感知到的是消毒水特有的清冽气味，然后是左臂和右腿传来的、被妥善包裹后依然沉闷而持续的痛感。789费力地睁开沉重的眼皮，视线花了片刻才适应室内明亮的光线。

白色的天花板，熟悉的校医院病房布置。他微微动了动脖子，目光向下移去。

第一眼看到的，是趴在床边的一个毛茸茸的脑袋。278就那样蜷缩在并不舒适的椅子上，上半身伏在病床边缘，睡着了。他眼下是浓重的乌青，头发乱糟糟地黏在额前，嘴角还隐约能看到一点干涸的血迹混着尘土——是昨晚挣扎呼救时留下的痕迹。他甚至没换下那身已经变得破破烂烂、沾满泥污血渍的作训服。

阳光透过窗户，落在他疲惫不堪的睡脸上，有一种近乎脆弱的安静。但即使是在睡梦中，他的眉头也微微拧着，一只手还无意识地、轻轻地搭在789没有受伤的右臂上，仿佛怕他消失一样。

789静静地看着他，心脏像是被一只无形的手紧紧攥住，又酸又胀。昨晚悬崖边的生死一刻、蚀骨的寒冷、还有自己脱口而出的那句誓言，潮水般涌回脑海。而比那些更清晰的，是278滚烫的眼泪和撕心裂肺的呼喊。

似乎是被789细微的动作惊动，278的眼睫颤了颤，猛地抬起头。四目相对。

278的眼中先是茫然，随即迅速被巨大的惊喜和仍未散去的后怕取代。“你醒了？！”他的声音沙哑得厉害，带着刚睡醒的混沌，却立刻伸手想去按呼叫铃，“医生！你等等，我叫医生！”

“别叫。”789开口，声音同样干涩低哑，“我没事。”

他的手动了动，下意识想阻止278，却牵动了左臂的伤，疼得他倒抽一口冷气。278立刻缩回按铃的手，转而紧张地扶住他完好的右臂，“你别乱动！医生说你就是太乱动才会失血过多晕过去！”

两人一时都沉默了。昨夜的惊心动魄和眼下病房的宁静形成了巨大反差，那些未说开的话、横亘在中间的裂痕，在此刻显得无比清晰。

278低着头，目光落在789被厚重石膏固定的左臂和右腿厚厚的绷带上，嘴唇翕动了几下，终于闷闷地开口，声音里带着浓重的鼻音和愧疚：“对不起……要不是为了救我，你也不会……”

“跟你没关系。”789打断他，语气很肯定，“是装备的问题。”他顿了顿，目光锐利地看向278，“倒是你，昨晚……吓傻了？”

278猛地抬头，眼圈瞬间又红了，像是被这句话戳中了最委屈的地方。他瞪着789，几乎是吼了出来，带着一种破罐破摔的坦白：“是！我是吓傻了！我他妈快吓死了！我当时满脑子都是你要是……你要是……我他妈……”

他哽住了，深吸了好几口气，才像是下定了决心，把憋在心里许久的话一股脑倒了出来，语速快得惊人：

“我接近C，是因为他每门课的笔记都总结得特别好！特别清楚！我……我是想抄过来，等期末的时候给你复习用的！你训练那么忙，根本没时间整理这些！我知道你因为成绩的事心里一直不痛快……”

他越说声音越低，最后几乎变成了嘟囔，带着点自暴自弃的沮丧：“谁知道你他妈会为这个生气……还跑去玩命加练……最后还搞成这样……”

空气再次安静下来。789彻底愣住了，他看着278因为激动和委屈而泛红的眼眶，看着他那副“千错万错都是我的错”的狼狈样子，心脏那片最柔软的地方被狠狠击中。

原来是这样。

那些他以为的亲密、那些让他烦闷不堪的醋意，源头竟然是这样笨拙、却纯粹到有点傻气的打算。是为了他。

所有的误解、猜忌和连日来的压抑，在这一刻，被278这番带着哭腔的坦白冲刷得干干净净，只剩下巨大的震动和铺天盖地的心疼。

“傻子……”789低声说，声音里带着他自己都未察觉的温柔和如释重负。

278没听清，或者说，他沉浸在自己的情绪里，没注意到789的变化。他的目光无意识地滑过789因为病号服宽松而微微敞开的领口，落在那片轮廓分明、随着呼吸轻轻起伏的胸腹区域。即使此刻伤痕累累地躺在病床上，789长期锻炼出的紧实肌肉线条依然清晰可见。

278的眼神里闪过一丝纯粹的、毫不掩饰的羡慕，还有更深层的、连他自己可能都未曾明晰的依赖与眷恋。他鬼使神差地伸出手，指尖带着一点点小心翼翼的颤抖，极其轻柔地触碰了一下789腹部那道因为呼吸而微微绷紧的肌肉沟壑。

触感温热而坚实，充满了生命的力量。这是将他从悬崖边拉回来的力量，是他在绝望中紧紧抓住的依靠。

“真好啊……”278喃喃自语，声音轻得像梦呓，充满了难以言喻的深情，“你可得……快点好起来。”

这个触碰短暂得如同错觉，却像一道电流，瞬间击穿了两人之间所有残余的隔阂。789的身体微微一僵，却没有躲闪，只是深深地望着278，看着他眼中那份混杂着羡慕、愧疚、依赖和某种更深沉情感的复杂光晕。

所有言语都成了多余。

789用没受伤的右手，缓缓地、坚定地覆上了278还停留在他腹肌上的那只手，将它紧紧握住。

“嗯。”他应了一声，承诺般说道，“会的。”

阳光洒满病房，尘埃在光柱中安静地飞舞。紧握的双手，无声地确认了比兄弟更亲密、比承诺更厚重的情感。裂痕在这一刻，被一种名为“真相”和“心疼”的粘合剂，彻底弥合，并且淬炼得更加坚固。
\newpage
\large
\chapter{风雨同舟}
\normalsize
\section{催化剂的考验}
\songti\zihao{4}
789受伤后的第三天，情况稳定了下来，但仍需住院观察。278几乎是长在了医院，课也不去上了，吃喝拉撒都在病房解决，眼睛底下的乌青就没消过，但精神头却好了很多，忙前忙后地给789削水果、倒水，嘴里絮絮叨叨，仿佛要把之前冷战没说的话都补回来。

下午，278被789强行赶回学校洗澡换衣服，顺便去辅导员那儿补个假条。他刚走没多久，病房门就被轻轻敲响了。

“请进。”789以为是护士。

门推开，进来的却是C。他手里提着一袋水果，脸上带着恰到好处的关切笑容。“789，听说你受伤了，好点了吗？”他走进来，把水果放在床头柜上，“278呢？刚在楼下好像看到他了。”

“回去换衣服了。”789言简意赅，他对C的观感很复杂。虽然误会解开了，知道278接近他是为了笔记，但那种因C而产生的、自身领域被侵入的不适感，并未完全消失。

C拉了把椅子在床边坐下，寒暄了几句伤势和集训的情况，话题却不知不觉地，引向了278。

“这次真是万幸，你没事，278也没大事。”C感叹着，随即话锋微妙地一转，语气温和，眼神却带着一种审视，“不过，经过这次……789，你有没有想过，你觉得……你的世界，真的对他最好吗？”

789的心猛地一沉，抬眼看向C。

C似乎没觉得自己的话有多尖锐，继续平静地分析，像在讨论一道习题：“我没有别的意思。只是，你看，你们在一起，好像总是……比较‘激烈’。”他斟酌着用词，“训练、比赛，现在又是这样差点出人命的意外。278那个人，看起来大大咧咧，其实心思没那么粗。他适合的， maybe是一种更……平稳一点的环境？”

他顿了顿，目光扫过789打着石膏的手臂和绷带缠绕的腿，声音放得更轻，却像针一样扎人：“你给他的这种生活，充满了汗水和风险，甚至……危险。这真的是他想要的，或者说，对他最好的选择吗？你有没有问过他，他会不会……害怕？”

这些话，像是一把钥匙，精准地捅开了789内心深处那扇从未敢轻易触碰的门。他一直以来的隐忧——自己是否能给278一个安稳的未来，自己所在的这个充满竞争和身体挑战的世界，是否会让278某天感到疲惫甚至后悔——被C如此直白地摊开在了阳光下。

789的脸色有些发白，嘴唇紧抿。他无法否认，C的每一个字，都敲打在他的软肋上。

C看着他的反应，轻轻叹了口气，像是无奈，又像是某种最后的确认。他忽然换了一种语气，带着点纯粹的、客观的欣赏，补充了一句，这句话却比之前的任何质疑都更具杀伤力：

“而且，说真的，278他……长得真的很帅，性格也好。他其实有很多选择，完全可以过上一种更轻松、更被大多数人认可的生活。你……确定要把他绑在你这条路上吗？”

病房里陷入了死一般的寂静。阳光透过窗户，照在789没有血色的脸上，他的拳头在身侧悄悄握紧，指甲陷进了掌心。

就在这时，病房门被“哐”一声推开。278顶着一头湿漉漉的头发，换了一身干净衣服，浑身散发着皂角的清新气息，像一阵风似的冲了进来。

“搞定！辅导员那边假条批了！”他话音落下，才察觉到病房里诡异的气氛。他看看脸色难看的789，又看看坐在一旁、表情平静无波的C，笑容僵在了脸上。

“C哥？你来了？”278的语气带着点不易察觉的警惕，他几乎是下意识地快步走到789床边，站的位置，刚好隔开了C和789的直接视线，形成一个保护的姿态。

C站起身，脸上又挂起了那副无可挑剔的温和笑容：“嗯，来看看789。你回来了就好，我也该走了，系里还有点事。”他对着789点点头，“好好休息，别多想。”

C离开后，病房里再次安静下来。278皱着眉，看向一言不发、眼神晦暗的789，心里咯噔一下。

“他……跟你说什么了？”278小心翼翼地问。

789抬起头，目光极其复杂地看向278。阳光勾勒出278刚洗过澡后格外清爽帅气的侧脸，C那句“他长得真的很帅”仿佛又在耳边响起。是啊，278有资本过上更轻松的生活，他凭什么……

“他问我，”789的声音干涩，几乎是一个字一个字地往外挤，“我的世界……是不是真的对你好。”

278愣住了，随即，一股怒火“噌”地窜上头顶。他猛地转头看向门口，仿佛C还没走远似的，气得胸口起伏：“他他妈什么意思？！”

“他说，”789闭了闭眼，像是用尽了全身力气，“我们在一起太‘激烈’，充满危险……他说你可能会害怕……说你值得更安稳的生活。”

“放他娘的屁！”278彻底炸了，他转回头，双手猛地抓住789没受伤的右肩，眼睛因为愤怒和委屈瞪得通红，“他懂个屁！什么是好？什么是安稳？跟他一样天天泡自习室、争保研名额就是好吗？！”

他死死盯着789的眼睛，声音因为激动而发颤，却带着一种前所未有的清晰和坚定：“789你给老子听好了！我278没那么怂！我要是怕，当初就不会进国护队！我要是想图安稳，我早就跟宿舍那帮人一样天天打游戏混日子了！”

他的手指用力，几乎要掐进789的肩膀肉里：“是，训练是苦，比赛是累，这次是差点没命！但我告诉你，我从来没后悔过！跟你在一块，流汗也好，流血也罢，老子心里踏实！什么叫死生契阔？就是他妈的在悬崖边上，你扑过来替我挡着的时候，我就知道了，这辈子就是你了！别人说的那种‘好’，那种‘安稳’，老子不稀罕！”

278喘着粗气，一番吼完，眼圈比刚才更红了，但眼神却亮得惊人，像烧着两团火。他看着被他一顿吼震住的789，语气忽然软了下来，带着点鼻音，委屈巴巴地说：

“你少他妈听别人瞎哔哔……你再敢因为这种屁话胡思乱想，我……我他妈就……”

他“我”了半天，也没想出什么有威慑力的威胁，最后只能恶狠狠地、带着哭腔总结：“我就不理你了！”

看着眼前这个像被踩了尾巴的猫一样炸毛、却又无比认真地向自己宣告“死生契阔”的278，789心中所有的阴霾、所有的自我怀疑，在那一刻，被这团炽热、笨拙却无比真诚的火焰，烧得干干净净。

他忽然低低地笑了起来，笑声从胸腔震动出来，牵动了伤口，让他一边笑一边抽气，可眼底的阴郁却一扫而空，重新变得明亮而坚定。

他抬起没受伤的右手，覆上278紧紧抓着自己肩膀的手，轻轻拍了拍。

“好。”他看着278，目光深沉如海，却清晰地倒映着对方的影子，“我知道了。”

C带来的那点阴云，被278这番毫无章法却情感磅礴的宣言，彻底驱散。他们的路或许崎岖，但此刻，两颗心前所未有地紧靠在了一起，风雨同舟。
\normalsize
\section{并肩为王}
\songti\zihao{4}
集训最终竞赛日在一种微妙的气氛中拉开序幕。各校队伍齐聚综合竞赛场，项目涵盖了障碍越野、战术协同、目标侦察与通信保障等多个环节。K大队伍的氛围尤其凝重，不仅因为之前与其他队伍的小摩擦，更因为他们的王牌——789，是吊着一只胳膊、拖着一条伤腿站在队列里的。

他的出现本身，就引来不少或惊讶、或同情、或隐含轻视的目光。H大那几个队员更是毫不掩饰地交换着眼神，仿佛在说“K大是没人了吗，让个重伤号上来充数？”

789对这一切视若无睹。他的脸色依旧有些苍白，但眼神却比以往任何时候都更加锐利和沉静，像一头蛰伏的、等待出击的猛兽。278站在他侧后方，不再是平时那副散漫样子，他仔细地最后一次检查了789身上的装备，特别是伤处的固定和缓冲，确保万无一失。他的眼神里没有了戏谑，只有全神贯注的谨慎。

“放心，撑得住。”789低声道。

“废话，我调整过的装备，绝对比你这原装的好用。”278嘴上不饶人，手下动作却轻柔无比。

竞赛开始。障碍越野环节，789的表现震惊了所有人。他无法像以前那样迅如闪电，但每一个动作都充满了计算好的精确和一种近乎残忍的坚忍。翻越高墙时，他主要依靠单臂和腰腹核心力量；匍匐通过低桩网时，他受伤的腿拖在后面，每一次摩擦都带来钻心的疼，但他哼都没哼一声，速度甚至不比完全健康的队员慢多少。那是一种用意志力对抗物理极限的表演。

278没有参赛资格，他作为后勤人员，守在场地边缘的补给和故障处理区。但他的注意力从未离开过赛场，尤其是789的身影。他手里拿着便携式通讯器，随时准备应对突发状况。

转折点发生在“目标侦察与信息回传”环节。各队需要深入模拟战场，利用携带的战场通讯设备，将侦察到的“敌情”加密传回指挥部。然而，就在K大小队即将抵达关键侦察点时，他们佩戴的通讯耳麦里，突然传来一阵极其尖锐、持续不断的频段干扰噪音！

“滋——！！！！！”

噪音瞬间刺穿了所有队员的耳膜，不仅无法接收指令，连基本的小队内部沟通都被切断。更糟糕的是，裁判组很快发现，这并非针对K大的恶意干扰，而是整个竞赛区域使用的公用频段，不知何故被一个未知的、强大的异常信号侵占了！其他队伍也陆续出现了通讯不畅的问题，比赛一度陷入停滞。

裁判和技术人员紧急排查，一时却找不到干扰源。时间一分一秒过去，赛场气氛变得焦躁。如果不能及时恢复通讯，整个竞赛项目可能作废，或者各队成绩都将大打折扣。

就在这混乱的时刻，278的身影动了。他没有像其他人一样茫然等待，而是迅速冲向K大的装备存放区，一把抓过备用的通讯中继设备和一个频谱分析仪。他蹲在地上，双手飞快地操作起来，眼神专注得像是在破解一个复杂的密码。

他的动作没有丝毫学院派的刻板，充满了野路子的直觉和效率。他没有试图去定位遥远的、难以捕捉的干扰源，而是反其道而行之。

“789！听到吗？能断断续续听到一点就敲两下麦克风！”278对着还能勉强传出一点破碎声音的通讯器吼道。

很快，耳麦里传来两声微弱的敲击。

“好！听着！我现在要你那边配合！把你设备的天线角度，向左偏转15度，对，然后向下压大概5度！然后保持这个姿势别动！”278语速极快，发出指令。他这是在利用战场环境和设备天线的指向性，尝试人为制造一个局部的、针对K小队通讯链路的“静默区”，来规避主流向上的干扰。

赛场边，前来观摩的A，目光不由自主地被278吸引。她看着那个平时在课堂上被她认为是“木头脑袋”的男生，此刻正用一种她无法完全理解的、基于大量实践经验和空间想象力的方式，快速而精准地调整着参数。他没有计算复杂的公式，更像是一个顶尖的机械师在凭手感调试一台精密的仪器。

“偏转了！现在干扰小了点，但信号还是很弱！”789断断续续的声音传来。

“收到！保持姿势！我再给你加个‘buff’！”278一边说着，一边从自己的万能工具包里掏出一个……锡纸包的能量棒？他快速将锡纸撕下，揉成一团，然后巧妙地缠绕在中继设备天线的特定部位。

这个看似儿戏的举动，却让频谱仪上代表K小队信号的波段，猛地清晰、稳定了起来！

“通了！K大通讯恢复！”裁判区有人大喊。

在整个赛场依旧被干扰笼罩的背景下，K大小队的通讯频道里，响起了789清晰、冷静的侦察汇报声。他们率先完成了信息回传！

那一刻，所有人的目光都聚焦在了那个蹲在地上、用锡纸和直觉解决了连裁判组都头疼的技术难题的278身上。他脸上没有得意，只是松了口气，抹了把额头的汗，然后第一时间通过耳麦问：“789，你那边怎么样？腿还能撑住吗？”

看台上，A微微张开了嘴，眼中充满了难以置信的震撼。她看着278，又望向远处赛场上，那个明明浑身是伤，却凭借着惊人意志力完美完成所有战术动作的789。

她忽然想起了高中时那个她暗恋的、有天赋但并非最顶尖的学长。那个学长也曾用一些看似“不规矩”但极其巧妙的方法，解决过棘手的竞赛题目。而此刻，278身上闪耀的那种解决实际问题的、跳跃性的实践智慧，和789身上所体现的、超越肉体痛苦的极致坚忍，这两种与她潜意识里一直欣赏的学术天才迥异的特质，以一种无比强烈和直观的方式，叠加在一起，冲击着她的认知。

她一直信奉的、建立在公式和逻辑之上的价值体系，在这一刻，被另一种截然不同、却同样强大、甚至在某些层面更显珍贵的价值所撼动。

她看着相互扶持、在逆境中硬生生杀出一条血路的789和278，第一次清晰地认识到：原来，优秀和强大，真的可以有如此不同的形态。 他们的世界，或许充满汗水和风险，但也孕育着在平静书斋里永远无法想象的坚韧、急智和那种……“死生契阔”的羁绊。

她轻轻呼出一口气，目光复杂，低不可闻地喃喃自语：“……厉害。”

这声感叹里，不再有轻视，而是带着一丝她自己都未曾察觉的……佩服。

最终，K大凭借在极端逆境下的出色表现，尤其是通讯恢复后的高效协同，逆风翻盘，赢得了联合集训的综合优胜奖。当789代表K大，一瘸一拐却脊背挺直地上台领取奖杯时，全场响起了热烈的掌声。那掌声，是给胜利者，更是给不屈的意志和绝境的智慧。

278在台下，看着聚光灯下的789，笑得比阳光还灿烂。他们再一次，并肩为王。
\normalsize
\section{归途与新程}
\songti\zihao{4}
集训的喧嚣与荣誉如同退潮的海水，渐渐消失在日常生活的沙滩之后。789和278回到了熟悉的K大校园，日子仿佛又回到了原来的轨道，但某些东西，已经悄然改变。

最大的变化发生在789身上。那次生死边缘的经历，像一次彻底的淬炼，将他心中因学业而起的焦躁与虚火一并烧尽了。他不再像过去那样，在训练场和自习室之间撕裂般地挣扎，也不再因为一次不理想的分数就陷入漫长的自我怀疑。

他开始用一种管理军事科目训练表的方式，来规划自己的学习。时间被分割成一块块，什么时间复习电路，什么时间做电磁学习题，什么时间进行恢复性训练，条理清晰，严格执行。他依然会为一道难题皱眉，但不再空泛地焦虑，而是沉下心来，一遍遍演算，甚至开始主动啃噬那些曾经让他望而生畏的理论基础。那是一种卸下了沉重包袱后的、真正意义上的“踏实”。

278的变化则更为外显。他依然会吐槽老师划的重点“反人类”，依然觉得厚厚的教材像催眠符，但宿舍里的游戏声确实减少了。他开始背着书包出现在图书馆和自习室，虽然时常坐不到半小时就开始抓耳挠腮，但至少，他人在那里。

更令人惊讶的是，某天下午，在图书馆的角落，278正对着一本《数字电路设计》愁眉苦脸，嘴里咬着笔帽，发出轻微的“咔哒”声。这时，一个细弱蚊蚋的声音在他旁边响起：

“那……那个……同学，你这里……卡诺圈的画法好像……有点问题。”

278吓了一跳，抬头看见一个瘦瘦小小的男生站在旁边，穿着洗得发白的格子衬衫，怀里抱着几本厚厚的物联网专业书，脸都快埋进衣领里了，耳朵尖通红。男生身上有股淡淡的、像是洗衣液混合着阳光的好闻味道。

“啊？哪里错了？”278来了精神，赶紧把书推过去。他认识这个男生，是同级的，好像叫276，物联网工程的，据说模电数电学得极好，但性格极其内向，几乎不跟人说话。

276像是被278的热情吓到，往后缩了缩，但还是伸出细白的手指，指尖点在278的草稿纸上，声音更小了：“这、这里……相邻项没、没画全……所以最简表达式不对……”他说话有点结巴，但逻辑异常清晰，三言两语就点出了关键。

“哦——！原来如此！”278恍然大悟，重重拍了下脑门，“谢了兄弟！你太牛了！”

276被他的大嗓门和动作惊得肩膀一颤，像只受惊的小兔子，小声说了句“不、不客气”，抱起书就想溜。

“哎别走别走！”278眼疾手快地拉住他（力道控制得极轻，怕吓着他），脸上堆起真诚（甚至有点谄媚）的笑容，“276同学是吧？久仰大名！那个……能不能再帮我看看这道模电的题？反馈这块我老是搞不懂……”

276僵在原地，走也不是，留也不是，脸更红了，但还是低着头，小声说：“……可、可以。”

于是，278开始了向276的“求学”之路。与向A请教时那种带着紧张和挫败感不同，面对社恐又柔软的276，278反而有种“大哥罩着小弟”般的自然。他会买好饮料零食“上供”，会在276讲解时夸张地表示“茅塞顿开”，虽然问题依然天马行空，但态度是前所未有的认真。而276，在278这种“傻大个”式的热情和崇拜下，似乎也慢慢放松了一点点，讲解时虽然依旧不敢看人眼睛，但语句流畅了许多。

这种细微的、向上的趋势，也落在了某些旁观者眼里。

一次《电磁场理论》课后，人群熙攘着涌出教室。A收拾好书本，正准备离开，一个身影有些犹豫地挡在了她面前。

是278。他脸上带着点不太好意思的笑，手里捏着本皱巴巴的笔记本。

“那个……A同学，能打扰你一分钟吗？”他的语气比之前任何一次请教都要诚恳。

A抬起头，平静地看着他，没说话，但也没有立刻走开，算是默许。

278赶紧翻开笔记本，指着一道关于电磁波极化的题目：“就是这道题，你课上讲的那种用复数表示法判断偏振态的方法，我大概听懂了，但我在想……如果是在我们之前集训遇到的那种复杂山地环境里，这种偏振的变化，会不会对那种突发性的频段干扰有什么影响？或者说，能不能利用这种特性来做点啥？”他的问题依然带着浓厚的“实践”色彩，甚至有些天马行空，但核心却紧紧扣住了课堂知识。

A的目光在278脸上停留了两秒。她看得出来，他不是在装模作样，也不是像以前那样只想套个答案。他的眼神里有一种新的东西——一种真正想要理解、并且试图将理论与他亲身经历的实际问题联系起来的求知欲。她甚至能感觉到，他最近应该经常和人讨论这类问题，思路比之前清晰了不少。

她沉默了几秒，然后接过他的笔记本，拿起笔，在一张空白页上快速画了几个坐标轴和向量。

“你这个想法，方向是可行的。”她的声音依旧清冷，但语气不再是纯粹的拒绝，“忽略山地复杂散射的初级模型下，偏振态的改变确实会影响信号接收质量。但你要考虑介质的各向异性……”

她开始讲解，语速依旧很快，逻辑跳跃，但这一次，她罕见地用了几个278能听懂的、贴近他知识背景的比喻，比如将偏振比作“信号的姿势”，干扰则是“错误的姿势导致了能量损耗”。

278听得极其认真，时不时提出一两个虽然不专业却切中要害的问题。这一次，A没有流露出不耐烦，反而会因为他的某个问题而稍作停顿，补充一两个关键点。

短暂的交流结束，278心满意足地抱着笔记本，连声道谢。A只是微微颔首，便转身离开。

走在回寝室的路上，A的思绪却有些飘远。她想起集训赛场上那个在干扰中精准“调频”的278，想起那个拖着伤腿却目光沉静如山的789。她想起刚才278提问时眼里那簇小小的、却真实存在的火苗。她甚至想起了最近偶尔在图书馆看到的，278和那个物联网系的、瘦小的男生凑在一起讨论问题的画面。

她看得出，278开始积极向上了。这种转变，或许笨拙，却带着一种从泥土里生长出来的、蓬勃的生命力。而这一切变化的源头，似乎都指向那个叫789的人，以及他们之间那种难以言喻的、能够彼此重塑的羁绊。

“还不赖。”

她望着远处操场上正在进行恢复训练的、那个熟悉的身影，用只有自己能听到的声音，轻轻说了一句。这一次，她的语气里，不再有居高临下的评判，反而带着一丝极淡的、对于某种顽强生命态的认可。

回归日常，并非回到原点。而是带着伤痕与荣耀，踏上了一段全新的征程。前路依旧有学业的大山，有未来的迷雾，但这一次，他们是真正的风雨同舟。
\newpage
\chapter{灯火归处}
\normalsize
\section{现实的重量}
\songti\zihao{4}
时光从不为任何人的青春驻足。转眼间，喧嚣躁动的大三下学期，便带着不容置疑的现实感，悄然降临。校园里开始弥漫起一种不同于以往的气息——招聘宣讲会的海报贴满了布告栏，考研自习室的位置需要提前一天抢占，同学们交谈的话题，不知不觉从游戏、八卦，转向了实习、简历和未来的去向。

温暖的春风拂过校园，吹开了新一轮的繁花，却吹不散悄然压在每个人心头的、名为“未来”的重量。789和278，也未能例外。

一天下午，789刚从一场面试回来。他穿着略显紧绷的西装，坐在宿舍书桌前，看着电脑屏幕上那封新收到的邮件，久久没有动作。邮件来自一家业内颇有名气的大型科技企业，内容简洁而有力：鉴于他在校期间出色的领导力、团队协作能力以及国护队的特殊经历，诚邀他毕业后加入公司的“行政管理培训生”项目。

这是一个难得的机会。稳定的薪资、清晰的晋升路径、体面的工作岗位。对很多毕业生来说，这堪称一份理想的offer。这无疑是对他过去三年在“另一条赛道”上拼搏价值的某种肯定。

但789的脸上，却看不到太多喜悦。他的目光掠过邮件末尾的公司地址——那是一座距离K大所在城市近千公里的沿海都市。

几乎在同一时间，278正接着一通来自家里的长途电话。电话那头，母亲的声音一如既往的温柔，却带着不容置疑的期盼：

“……小辉啊，你爸托人问过了，咱们市里的供电公司今年有名额，专业对口，稳定离家又近。你大学也玩……也锻炼三年了，该收收心回来了。你一个人在外面，爸妈总是不放心……”

278靠在阳台栏杆上，望着楼下熙攘的人群，含混地“嗯嗯啊啊”应着。挂断电话后，他烦躁地抓了把头发。回乡，进入一个安稳的体制内单位，过着父母眼中“正常”的生活——这本该是顺理成章的选择，可此刻，他却感到一种莫名的抗拒和……窒息。

他回到寝室，看到789对着电脑屏幕出神的背影，那身还没来得及换下的西装勾勒出他宽厚的肩背，却莫名透着一丝紧绷。278心里咯噔一下，一种不祥的预感浮现。

“喂，面试咋样？”他故作轻松地凑过去，手臂习惯性地搭上789的肩膀。

789身体几不可察地僵了一下，随即若无其事地最小化了邮件窗口，淡淡道：“还行。”

“有戏？”278追问，仔细打量着789的侧脸。

“嗯。”789应了一声，站起身，开始解西装的扣子，动作有些缓慢，“给了个管培生的offer。”

“卧槽！牛逼啊789！”278眼睛一亮，由衷地为他高兴，用力拍了下他的背，“我就知道你可以！请客！必须请客！”

789扯了扯嘴角，算是回应，但笑容并未到达眼底。他脱下西装外套，小心地挂好，仿佛那是什么易碎品。

傍晚，两人像过去无数个日子一样，并肩在操场散步。夕阳将他们的影子拉得很长。但这一次，气氛却不同以往。一种沉重的东西横亘在两人之间，沉默像不断膨胀的实体，挤压着熟悉的默契。

最终还是278先憋不住了，他踢着跑道上的小石子，闷闷地开口：“我妈……想让我回去。说家里给找了工作。”

789的脚步几不可察地顿了一下，没有看278，只是“嗯”了一声，表示知道了。

又是一阵令人窒息的沉默。

“你呢？”278停下脚步，转过身，直视着789，眼神里带着他自己都没察觉的紧张和期盼，“那个offer……在哪？”

789也停了下来，他终于看向278，目光复杂得像一团纠缠的线。夕阳的余晖落在他眼里，明明灭灭。

“深圳。”他吐出两个字。

仿佛一道无声的惊雷在两人之间炸开。

深圳。千里之外。一个意味着分离、未知和巨大变数的地名。

278脸上的最后一点强装的笑意也消失了。他张了张嘴，想说什么，却发不出任何声音。所有关于未来的想象，似乎都在这一刻，被现实无情地撕开了一道巨大的裂口。

789看着278瞬间黯淡下去的眼神，看着他脸上那毫不掩饰的失落和无措，心脏像是被一只冰冷的手狠狠攥住。他想起C曾经那个尖锐的问题：“你觉得你的世界，真的对他最好吗？”

此刻，这个问题带着千钧的重量，再次压上他的心头。他选择的这条路，他可能要去往的远方，真的是278想要的吗？还是……会成为278的负担和枷锁？

现实的重量，从未如此刻般清晰而冰冷地，落在了他们年轻却已然紧密相连的肩头。灯火可期的归处，似乎在这一刻，变得模糊而遥远。

\normalsize
\section{我的选择是你}
\songti\zihao{4}
那层窗户纸被捅破后，压抑的气氛持续了好几天。278变得异常安静，不再咋咋呼呼地拉着789讨论未来，也不再抱怨功课。他只是沉默地陪着789去图书馆，沉默地一起吃饭，沉默地在夜晚的操场上一圈圈走着，但两人之间的空气像是凝固了，以往那种无需言说的默契被一种小心翼翼的试探和不安取代。

789同样备受煎熬。深圳的那个offer像一块烫手的山芋，他几次点开邮件，指尖悬在“回复”键上，却迟迟无法落下。每一次，他眼前都会闪过278听到“深圳”两个字时瞬间黯淡的眼神。他第一次对自己坚信不疑的道路产生了动摇。他的世界，他的拼搏，如果最终通往的是与这个人的分离，那这一切的意义又在哪里？

转折发生在一个深夜。278破天荒地没有像往常一样在789的宿舍磨蹭到熄灯，而是早早说了句“累了，先回去了”，便离开了。789独自一人坐在书桌前，对着电脑屏幕上空白的回复邮件界面，只觉得心烦意乱。他索性关掉电脑，拿起烟盒和打火机，走到了空旷的阳台。

夜凉如水，远处城市的灯火像撒在地上的碎星。他刚点燃一支烟，就听到身后传来熟悉的脚步声。一回头，看见278穿着背心裤衩，踩着拖鞋，也走上了阳台，手里居然也拿着一罐啤酒。

“睡不着？”789哑着嗓子问。

“嗯。”278走到他身边，靠在栏杆上，拉开啤酒罐，仰头灌了一大口。冰凉的液体划过喉咙，他舒服地叹了口气，然后陷入沉默。

两人并排站着，望着远处的灯火，谁都没有先开口。烟头的红点在789指间明灭，啤酒罐壁凝结的水珠沾湿了278的手指。这种沉默不同于之前的压抑，更像是一种风暴来临前的平静，有什么沉重的东西正在积蓄，等待破土。

终于，278像是下定了决心，将剩下的啤酒一饮而尽，空罐子被他捏得咔咔作响。他转过身，面对789，夜色中，他的眼睛亮得惊人，带着一种破釜沉舟的坚决。

“789，”他开口，声音因为酒精和紧张而有些沙哑，却异常清晰，“我跟我妈吵了一架。”

789夹着烟的手指顿住了，抬眼看他。

“我告诉她，我不回去。”278一字一顿地说，像是在宣布一个不容更改的决定，“那个供电局的工作，谁爱去谁去。我278，不伺候。”

他往前凑近一步，几乎要贴上789，仰着头，目光灼灼地锁住789的眼睛：“我知道你在想什么。你觉得你的路太苦，你的世界太险，怕我跟不上，怕我后悔，是不是？”

789喉咙发紧，想否认，却发不出声音。278的眼神太锐利，直接剖开了他内心最深处的隐忧。

“我告诉你，789，”278的声音陡然提高，带着一种被误解的委屈和愤怒，“你少他妈自作聪明！老子这辈子做过最不后悔的事，就是跟你进了国护队，就是那天在悬崖边上没撒开你的手！”

他深吸一口气，像是用尽了全身的力气，说出了在心底盘旋了无数遍的决定：

“所以，你也别想甩开我。听着，你去深圳，老子就跟你去深圳！你去天涯海角，老子也奉陪到底！你去哪里，我的后勤保障就跟到哪里！这话我说的，一辈子算数！”

这番话，像一阵狂风，瞬间吹散了789心中所有的迷雾和挣扎。他看着眼前这个眼眶泛红、像头被逼到墙角却更加凶狠的小兽般的278，看着他因为激动而微微颤抖的身体，一股滚烫的热流猛地冲垮了他所有的理智和顾虑。

他扔掉烟头，猛地伸出手，一把将278紧紧地、用力地箍进了怀里。力道之大，像是要把他揉进自己的骨血里。

278被他勒得差点喘不过气，但只是僵硬了一瞬，便立刻反手更紧地抱住了789的腰，把脸深深埋进他带着汗味和烟味的颈窝里。

“笨蛋……”789的声音沙哑得厉害，带着不易察觉的颤抖，滚烫的呼吸喷在278的耳廓，“……那地方，消费高，竞争强……活不好干。”

“怕个毛！”278闷在他怀里，声音嗡嗡的，却斩钉截铁，“有老子在，还能饿着你？保证把你养得白白胖胖！”

789低低地笑了起来，胸腔震动，连带着怀里的278也一起振动。他收紧了手臂，低下头，嘴唇贴着278的短发，用一种前所未有的、郑重如誓言般的语气，沉声说：

“好。那以后……我的后背，就只交给你了。有你在身后，我哪里都敢闯。”

没有华丽的辞藻，没有浪漫的承诺。只有最朴素的交付，和最坚定的追随。在这个平凡的深夜里，两个年轻人用他们特有的方式，完成了对彼此、也是对未来的最终选择。

现实的重量依然存在，前路的挑战只多不少。但此刻，紧紧相拥的他们，仿佛拥有了对抗整个世界的勇气。

我的征程，是你。

我的归处，也是你。
\normalsize
\section{卒业式}
\songti\zihao{4}
盛夏的阳光毫无保留地倾泻在K大的体育场，毕业典礼的会场充满了喧闹、鲜花和离别的笙箫。空气中弥漫着栀子花的香气，混合着青春的感伤与对未来的憧憬。

789和278穿着一样的学士服，坐在同一片方阵里。宽大的袍子暂时掩盖了789挺拔的身姿和278匀称的骨架，让他们看起来和周围每一个毕业生并无不同。但当你仔细看他们的眼睛——789眼神沉静坚定，278眼中带着一如既往的明亮和一丝不易察觉的紧张——便会知道，他们共同走过的这四年，是何等的波澜壮阔。

校长致辞，拨穗正冠。当学院院长将789学士帽上的流苏从右边拨到左边，微笑着与他握手祝贺时，789的心跳平稳而有力。他接过那卷象征着学业的证书，却没有立刻走下台。

在全场数千名师生的注视下，他做了一个出人意料的举动。

他转过身，目光精准地投向台下那个一直紧紧望着他的方向。然后，在278惊愕的目光中，在渐渐响起的窃窃私语声中，789一步一步，坚定地走到了278的面前。

278完全懵了，呆呆地看着789，忘了反应。

789没有说话，只是深深地望进278的眼睛，那里面有紧张，有疑惑，更有全然的信任。他抬起手，将自己头顶那顶刚刚被拨穗的、象征着圆满与认可的黑色学士帽，轻轻取了下来。

接着，在所有人——包括台上的院长、台下的教授和同学——或惊讶、或好奇、或了然的目光中，789伸出双手，无比郑重地、带着一种近乎仪式般的庄重，将属于自己的学士帽，端端正正地戴在了278的头上。

流苏垂在278的额前，微微晃动。

那一刻，时间仿佛静止。悬崖边的寒风、黑夜里的誓言、病房中的相守、深夜阳台的抉择……四年的光阴与生死，都凝聚在这一戴之中。

“答应了你的。”789的声音不高，却清晰地传入278耳中，带着如释重负的温柔和一丝不易察觉的哽咽，“给你戴上学士帽。”

278的视线瞬间就模糊了。眼泪毫无征兆地涌出，冲垮了他所有的故作镇定。他什么都说不出来，只是重重点头，泪水砸在崭新的学士服上，洇开深色的水渍。他猛地伸手，紧紧抓住了789的手腕，力道大得指节发白。

周围安静了一瞬，随即，爆发出了雷鸣般的、经久不息的掌声。这掌声里，有理解，有感动，更有对这份超越常规的深厚情谊的深深祝福。

当晚，喧嚣散尽。两人换下了学士服，穿着简单的T恤短裤，回到了他们最熟悉的地方——那片洒满了无数汗水和秘密的操场。

夜空如洗，星子寥落，晚风温柔。跑道上是夜跑的学生，看台上是窃窃私语的情侣，一切都和过去的无数个夜晚一样。

789和278并肩走到操场中央，这里曾是他们加练、争吵、也是和解和确立心意的地方。

789停下脚步，从裤兜里掏出一个小巧的丝绒盒子。打开，里面是两枚款式简单却闪着温润光泽的素圈戒指。

他拿起稍小的一枚，在278震惊的目光中，什么也没问，只是执起他的左手，缓缓地、坚定地，将戒指套进了他的无名指。尺寸刚刚好。

然后，他将另一枚戒指和盒子递给278。

278的手抖得厉害，试了好几次，才终于将另一枚戒指，同样郑重地戴在了789左手的无名指上。

没有单膝跪地，没有海誓山盟。只有星光为证，清风为媒。

789抬起手，看着指间那圈微光，然后紧紧握住了278同样戴着戒指的手。两枚素圈在月光下轻轻相碰，发出微不可闻的脆响。

“差不多了。”789说。

“嗯。”278点头，用力回握。

镜头拉远，两个曾行走在逆光中的少年，身影在广阔的操场上逐渐变小，却始终双手紧握，坚定不移地向着灯火通明的校园外走去，携手踏上属于他们的光明坦途。

死生契阔，与子成说。执子之手，与子偕老。
\newpage
\chapter*{后记：交织的善意}
\normalsize
\songti\zihao{4}
毕业典礼散场的人潮熙攘攘。早已拿到T大直博offer的A，抱着自己的证书，正随着人流往外走。她依旧是一副冷静淡然的样子，与周围的兴奋和不舍格格不入。

忽然，她的目光捕捉到了两个熟悉的身影——是278和789，他们正被人围着道贺，278笑得见牙不见眼，789站在他身边，眼神温和。278的学士服口袋里，露出了一个科技公司实习邀请函的角落。

人潮涌动间，A与278擦肩而过。就在那一瞬间，A的脚步几不可察地顿了一下。

她停了下来，转过身，罕见地主动开口，叫住了即将走过去的278：

“278。”

她的声音依旧平淡，没有什么起伏。

278和789都诧异地回过头。

A看着278，目光在他脸上停留了一秒，仿佛透过他，看到了那个在集训赛场上用锡纸和直觉解决故障的身影。她微微颔首，语气简洁：

“祝你顺遂。”

然后，在278惊愕的目光中，她像是完成一个通知程序般，补充了几句，声音不大，却清晰无误：

“你去实习的那家公司，技术部的总监，是276的叔叔。”

她顿了顿，仿佛在陈述一个客观事实：“276和他叔叔讲过你解决频段干扰的事。他们欣赏有急智的人。”

最后，她似乎只是随口一提，不带任何情绪地补充了另一个消息：“276已经去A大了。”

——那座拥有她高中时曾暗自欣赏过的、那位“并非第一梯队”学长的大学。这句话轻飘飘的，却像一把小小的钥匙，轻轻旋开了她心中某个尘封的角落，让一点光透了进去，无关风月，只为释然。

说完，她不看278瞬间瞪大的眼睛和789若有所思的表情，甚至没有等任何回应，便干脆利落地转身，抱着她的证书，汇入人海，背影清瘦而决绝，很快消失不见。

这个举动，融合了她和276两人所能表达的、最极致的善意。一种基于理性认可，却超越了常规交际的、笨拙而珍贵的祝福。

278站在原地，愣了很久，才喃喃道：“276他……还有A她……”

789看着A消失的方向，又看了看指间的戒指，最后将目光落回还在懵懂的278脸上，嘴角缓缓勾起一个深刻的弧度。这个世界，似乎比他们想象的，要温柔那么一点点。

“嗯。”他握紧了278的手，“走了，回家。”
\end{document}